\chapter{Introduction}
\label{ch:Introduction}%

\section{Purpose}
\label{sec:Purpose}%

\par The purpose of this Design Document (DD) is to provide a detailed and structured
description of the architectural and design choices adopted for the Best Bike Paths (BBP)
system. The document acts as a technical blueprint that guides the development team
through the implementation phase, clarifying how the system requirements identified in
the RASD are concretely addressed.

\par This document aims to ensure a shared and precise understanding of the system’s
technical solution among all stakeholders involved in the project. By explicitly describing
components, interactions, and design rationales, it helps align expectations, reduce
ambiguities, and support coordinated development activities.

\par Furthermore, the DD motivates and documents the main architectural decisions,
explaining the selection of technologies, architectural styles, and design patterns.
This information is intended to preserve design consistency over time and to serve as a
reference for future maintenance, evolution, or extension of the BBP system.

\section{Scope}
\label{sec:Scope}%

\par As described in the Requirements Analysis and Specification Document (RASD),
Best Bike Paths (BBP) is a software system designed to support cyclists in discovering,
recording, and evaluating bike paths based on their quality, safety, and overall suitability.
The system enables users to record cycling trips, report obstacles and street conditions,
and access aggregated information about routes shared by the community.

\par BBP integrates data collected manually by users and automatically through mobile
device sensors, enriching it with external information such as maps and weather conditions.
The system processes these inputs to compute scores and consolidated evaluations for bike
paths, supporting informed route selection for both registered and unregistered users.

\par The scope of this Design Document is limited to the technical realization of the BBP
system. It focuses on software architecture, component design, data management,
interaction flows, and deployment aspects, without redefining functional or non-functional
requirements already specified in the RASD.

\section{Definitions, Acronyms, and Abbreviations}
\label{sec:definitions-acronyms-abbreviations}%

\par This section reports the definitions, acronyms, and abbreviations used throughout
the document. The terminology is consistent with the RASD to ensure traceability and
conceptual coherence between requirements and design artifacts.

\section{Revision History}
\label{sec:revision-history}%

\begin{itemize}
      \item \textbf{Version 1.0} (2025-01-05): Initial release of the Design Document.
\end{itemize}

\section{Reference Documents}
\label{sec:reference-documents}%

\begin{itemize}
      \item Best Bike Paths (BBP) Requirements Analysis and Specification Document (RASD).
      \item R\&DD Assignment description, Academic Year 2025--2026.
      \item Provided course materials and lecture notes.
\end{itemize}

\section{Document Structure}
\label{sec:document-structure}%

\begin{enumerate}
      \item \textbf{Introduction}: Presents the purpose, scope, reference documents, and
            overall organization of the Design Document, providing the necessary context
            for understanding the subsequent sections.
      \item \textbf{Architectural Design}: Describes the high-level architecture of BBP,
            including major components, their responsibilities, interactions, and
            deployment environment. Architectural styles and key design decisions are
            discussed and justified.
      \item \textbf{Detailed Design}: Provides a deeper view of the internal structure of
            the system, including component interfaces, data models, and main design
            patterns adopted to satisfy the requirements.
      \item \textbf{User Interface Design}: Illustrates the design of the user interfaces,
            describing how cyclists interact with the system through mobile and, where
            applicable, web interfaces.
      \item \textbf{Requirements Traceability}: Maps the requirements defined in the RASD
            to the architectural and design elements described in this document, ensuring
            completeness and consistency.
      \item \textbf{Implementation, Integration, and Test Plan}: Outlines the planned
            development process, integration strategy, and testing activities required to
            validate the system.
      \item \textbf{Effort Spent}: Reports the workload distribution and effort spent by
            each member of the project team.
      \item \textbf{References}: Collects all bibliographic and technical references cited
            throughout the document.
\end{enumerate}
